\documentclass[11pt]{article}
\usepackage[a4paper,margin=1in]{geometry}
\usepackage{amsmath,amssymb,mathtools}
\usepackage{booktabs,array}
\usepackage{enumitem}
\usepackage[dvipsnames]{xcolor}
\usepackage{hyperref}
\usepackage{microtype}
\usepackage{comment}

\hypersetup{colorlinks=true,linkcolor=MidnightBlue,urlcolor=MidnightBlue,
  pdftitle={Recitation 5: Conditional Mean and Average Treatment Effect},
  pdfauthor={Xuanxi Zhang}}

% Default: student handout. Define \WITHSOLUTIONS before inputting this file
% to include answers and the instructor's 75-minute pacing guide.
\newif\ifsolutions
\ifdefined\WITHSOLUTIONS
  \solutionstrue
\else
  \solutionsfalse
\fi

\newcommand{\Pbb}{\mathbb{P}}
\newcommand{\E}{\mathbb{E}}
\newcommand{\rv}[1]{\widetilde{#1}}
\newcommand{\po}[1]{\widetilde{\mathrm{po}}_{#1}}
\newcommand{\indep}{\mathrel{\perp\!\!\!\perp}}
\newcommand{\ATE}{\operatorname{ATE}}

\ifsolutions
  \newenvironment{answer}{%
    \par\medskip\noindent\textcolor{MidnightBlue}{\textbf{Solution.}}\ }
    {\par\medskip}
  \newcommand{\workingspace}[1]{}
\else
  \excludecomment{answer}
  \newcommand{\workingspace}[1]{\vspace{#1}}
\fi

\newcommand{\keybox}[1]{%
  \begin{center}
  \fcolorbox{MidnightBlue}{MidnightBlue!5}{%
    \begin{minipage}{0.91\textwidth}#1\end{minipage}}
  \end{center}}
\setlist[enumerate]{label=(\alph*),leftmargin=2em,itemsep=0.4em}

\title{\huge Recitation 5\\[0.35em]
  \normalsize Conditional Mean and Average Treatment Effect\\[0.35em]
  \large\ifsolutions Solutions and Teaching Notes\else Review and Practice Problems\fi}
\author{Xuanxi Zhang}
\date{October 9, 2026}

\begin{document}
\maketitle
\noindent\textit{Reading:}
\href{https://github.com/cfgranda/ps4ds/blob/main/ps4ds_preprint.pdf}
{\emph{Probability and Statistics for Data Science}}, Sections 7.8--7.9.

\section{Conditional mean}
Fix an observed value $x$, then average $\rv{y}$ using its conditional
distribution. Assuming the relevant means exist and are finite,
\[
  \mu_{\rv{y}\mid\rv{x}}(x)
  :=\E[\rv{y}\mid\rv{x}=x]
  =\begin{cases}
    \displaystyle\sum_y y\,p_{\rv{y}\mid\rv{x}}(y\mid x),
       & \text{discrete }\rv{y},\\[0.9em]
    \displaystyle\int_{-\infty}^{\infty}y\,
      f_{\rv{y}\mid\rv{x}}(y\mid x)\,dy,
       & \text{continuous }\rv{y}.
  \end{cases}
\]
When both variables are discrete, normalize by
$p_{\rv{x}}(x)>0$:
$p_{\rv{y}\mid\rv{x}}(y\mid x)
=p_{\rv{x},\rv{y}}(x,y)/p_{\rv{x}}(x)$.

\keybox{\textbf{A function versus a random variable.}
$\mu_{\rv{y}\mid\rv{x}}(x)$ is a number once $x$ is fixed.
The \emph{conditional mean function} maps $x$ to that number.
Plugging in the random input gives the \emph{conditional mean}
$\mu_{\rv{y}\mid\rv{x}}(\rv{x})=\E[\rv{y}\mid\rv{x}]$,
which is a random variable.}

\paragraph{Iterated expectation (the tower law).}
Average the conditional means using the probabilities of the groups:
\[
  \boxed{\E[\rv{y}]=\E\!\left[\mu_{\rv{y}\mid\rv{x}}(\rv{x})\right].}
\]
For discrete $\rv{x}$ this becomes
\[
  \E[\rv{y}]=\sum_x p_{\rv{x}}(x)\,
    \mu_{\rv{y}\mid\rv{x}}(x).
\]
For continuous $\rv{x}$, replace the sum by an integral against
$f_{\rv{x}}(x)$. Group means get equal weight only when group probabilities
are equal.

\paragraph{From data.}
For observed pairs $X=\{(x_i,y_i):i=1,\ldots,n\}$, let
$I_x=\{i:x_i=x\}$. When $|I_x|>0$, estimate the conditional mean by
the within-group average $|I_x|^{-1}\sum_{i\in I_x}y_i$.
For a continuous predictor, exact repeats may be unavailable; one can average
nearby observations or fit a model for the conditional mean.

\newpage
\section{Average treatment effect}
Let $\rv{t}\in\{0,1\}$ denote treatment. Each unit has potential outcomes
$\po{0}$ under control and $\po{1}$ under treatment. The observed outcome
satisfies \textbf{consistency}:
\[
  \rv{y}=(1-\rv{t})\po{0}+\rv{t}\po{1}.
\]
Only the potential outcome for the received treatment is observed. The
population \textbf{average treatment effect} is
\[
  \boxed{\ATE=\E[\po{1}]-\E[\po{0}]
       =\E[\po{1}-\po{0}].}
\]

\paragraph{Observed contrast and randomization.}
The difference visible in the observed groups is
\[
  \Delta_{\mathrm{obs}}
  :=\mu_{\rv{y}\mid\rv{t}}(1)-\mu_{\rv{y}\mid\rv{t}}(0).
\]
In general this compares different populations. If
$\po{r}\indep\rv{t}$ for $r=0,1$, and both treatment groups have positive
probability, then
\[
  \mu_{\rv{y}\mid\rv{t}}(r)
  =\E[\po{r}\mid\rv{t}=r]
  =\E[\po{r}],
  \qquad \Delta_{\mathrm{obs}}=\ATE.
\]
Random assignment provides this independence. Differences of \emph{sample}
means still have sampling error; they estimate the population contrast.

\paragraph{Adjustment for a discrete confounder.}
Let $\rv{c}$ be a pre-treatment characteristic, such as smoking status.
Write $\mu_{\rv{y}\mid\rv{c},\rv{t}}(c,r)
=\E[\rv{y}\mid\rv{c}=c,\rv{t}=r]$. Assume:
\begin{enumerate}[label=\arabic*.,itemsep=0.2em]
  \item Consistency as above.
  \item Conditional independence:
    $\po{r}\indep\rv{t}\mid\rv{c}$ for $r=0,1$.
  \item Positivity: $0<\Pbb(\rv{t}=1\mid\rv{c}=c)<1$
    for every $c$ with $p_{\rv{c}}(c)>0$.
\end{enumerate}
Iterated expectation gives the adjustment formula:
\[
  \boxed{\ATE=\sum_c p_{\rv{c}}(c)
    \left[\mu_{\rv{y}\mid\rv{c},\rv{t}}(c,1)
          -\mu_{\rv{y}\mid\rv{c},\rv{t}}(c,0)\right].}
\]
Indeed, for either $r$,
\[
  \E[\po{r}]
  =\sum_c p_{\rv{c}}(c)\,\E[\po{r}\mid\rv{c}=c]
  =\sum_c p_{\rv{c}}(c)\,\mu_{\rv{y}\mid\rv{c},\rv{t}}(c,r).
\]
The second equality uses conditional independence and consistency.

\keybox{\textbf{Choose the weights before calculating.}
Adjustment uses the \emph{same target-population weights}
$p_{\rv{c}}(c)$ for both treatments. The raw group mean instead uses
$p_{\rv{c}\mid\rv{t}}(c\mid r)$, which may differ across treatment groups.
For a target consisting of all study participants, estimate the common weights
using the pooled counts.}

\newpage
\section*{Practice problems}
\noindent Four core problems; Problem 5 is optional. The indicated times
are for independent work before discussion.

\subsection*{Problem 1: Life expectancy (Exercise 7.11) \hfill [5 minutes]}
In the model in Exercise 7.11, age at death has mean $37.5$ years.
First-year deaths have probability $0.25$ and are recorded as age $0$.
Let $\rv{a}$ denote the recorded age at death, and let $\rv{s}$ equal $1$
if a person survives the first year and $0$ otherwise.
\begin{enumerate}
  \item Write $\E[\rv{a}]$ as a weighted average of
    $\mu_{\rv{a}\mid\rv{s}}(0)$ and
    $\mu_{\rv{a}\mid\rv{s}}(1)$.
  \item Find the mean age at death among people who survive the first year.
    Explain why this differs from the mean for the whole population.
\end{enumerate}
\workingspace{7.5cm}

\begin{answer}
First-year deaths are coded as zero, so
\[
  \mu_{\rv{a}\mid\rv{s}}(0)=0,
  \qquad \Pbb(\rv{s}=1)=1-0.25=0.75.
\]
By iterated expectation,
\begin{align*}
  \E[\rv{a}]
  &=\Pbb(\rv{s}=0)\mu_{\rv{a}\mid\rv{s}}(0)
    +\Pbb(\rv{s}=1)\mu_{\rv{a}\mid\rv{s}}(1),\\
  37.5&=0.25\cdot0+0.75\,\mu_{\rv{a}\mid\rv{s}}(1).
\end{align*}
Thus
\[
  \boxed{\mu_{\rv{a}\mid\rv{s}}(1)=\frac{37.5}{0.75}=50\text{ years}.}
\]
Conditioning removes the first-year deaths that lower the overall average.
The answer is an \emph{age at death}; it is not the expected number of
additional years after the first birthday, which would be $50-1=49$.

\paragraph{Discussion prompt.}
Ask students to identify the two group weights before doing arithmetic.
The statement $\E[\rv{a}]=37.5$ does not say that most adults died near age
$37.5$; a population mean alone does not describe the distribution of adult
lifespans. All calculations use the simplified model in the exercise.
\end{answer}

\newpage
\subsection*{Problem 2: A conditional mean from a model \hfill [8 minutes]}
\noindent\textit{Adapted from the uniform--Bernoulli exercise in the
Fall 2025 archive, \texttt{mine/r4}.}

\medskip
Let $\rv{u}\sim\operatorname{Uniform}(0,1)$, and conditional on
$\rv{u}=u$, let $\rv{x}\sim\operatorname{Bernoulli}(u)$.
Here $\rv{u}$ is explicitly random.
\begin{enumerate}
  \item Find $\Pbb(\rv{x}=1)$ and the conditional densities
    $f_{\rv{u}\mid\rv{x}}(u\mid1)$ and
    $f_{\rv{u}\mid\rv{x}}(u\mid0)$.
  \item Compute the conditional mean function
    $g(x):=\mu_{\rv{u}\mid\rv{x}}(x)$ for $x\in\{0,1\}$.
  \item Give the possible values and probabilities of the random variable
    $g(\rv{x})$. Verify that $\E[g(\rv{x})]=\E[\rv{u}]$.
\end{enumerate}
\workingspace{8cm}

\begin{answer}
\textbf{(a)} Integrating over the latent value $u$,
\[
  \Pbb(\rv{x}=1)=\int_0^1 u\,du=\frac12,
  \qquad \Pbb(\rv{x}=0)=\frac12.
\]
Bayes' rule for a continuous latent variable and a discrete observation gives
\[
  f_{\rv{u}\mid\rv{x}}(u\mid x)
  =\frac{p_{\rv{x}\mid\rv{u}}(x\mid u)f_{\rv{u}}(u)}
         {p_{\rv{x}}(x)}.
\]
Therefore, for $0<u<1$,
\[
  \boxed{f_{\rv{u}\mid\rv{x}}(u\mid1)=2u,
  \qquad f_{\rv{u}\mid\rv{x}}(u\mid0)=2(1-u),}
\]
and both densities are zero outside $(0,1)$.

\textbf{(b)} Average $u$ under each conditional density:
\[
  g(1)=\int_0^1u(2u)\,du=\frac23,
  \qquad
  g(0)=\int_0^1u\,2(1-u)\,du=\frac13.
\]
Hence $\boxed{g(x)=(1+x)/3}$ for $x\in\{0,1\}$.

\textbf{(c)} The random conditional mean is
$g(\rv{x})=(1+\rv{x})/3$. Its pmf is
\[
  \Pbb\!\left(g(\rv{x})=\frac13\right)=\frac12,
  \qquad
  \Pbb\!\left(g(\rv{x})=\frac23\right)=\frac12.
\]
Consequently,
\[
  \E[g(\rv{x})]
  =\frac12\cdot\frac13+\frac12\cdot\frac23
  =\boxed{\frac12}=\E[\rv{u}].
\]
Seeing $\rv{x}=1$ shifts the conditional mean upward; seeing $\rv{x}=0$
shifts it downward. Averaging over both possible observations recovers the
unconditional mean.
\end{answer}

\newpage
\subsection*{Problem 3: Coffee and life expectancy \hfill [10 minutes]}
\noindent\textit{Exercise 7.17.} The observational study gives these counts
and average lifespans:
\begin{center}
\begin{tabular}{@{}lrrrr@{}}
\toprule
& \multicolumn{2}{c}{Number of people}
& \multicolumn{2}{c}{Mean lifespan (years)}\\
\cmidrule(lr){2-3}\cmidrule(l){4-5}
& Coffee & No coffee & Coffee & No coffee\\
\midrule
Smokers     & 80 & 20 & 68 & 70\\
Non-smokers & 40 & 60 & 82 & 80\\
\bottomrule
\end{tabular}
\end{center}
Let $\rv{t}=1$ mean drinking coffee, $\rv{c}=1$ mean smoking, and
$\rv{y}$ denote lifespan. Use all study participants as the target population.
\begin{enumerate}
  \item Compute the unadjusted difference in sample means, coffee minus
    no coffee (the exercise's ``observed ATE'').
  \item Estimate the ATE after adjusting for smoking. Assume
    $\po{r}\indep\rv{t}\mid\rv{c}$ for $r=0,1$, as well as consistency
    and positivity. Show the weights you use.
  \item Explain the change after adjustment. Compare the two within-group
    differences with the adjusted estimate.
\end{enumerate}
\workingspace{6.5cm}

\begin{answer}
\textbf{(a)} The two observed treatment-group means are
\[
  \bar y_{\mathrm{coffee}}=\frac{80(68)+40(82)}{120}
       =\frac{218}{3},
  \qquad
  \bar y_{\mathrm{no\ coffee}}=\frac{20(70)+60(80)}{80}
       =\frac{155}{2}.
\]
Thus the unadjusted contrast is
\[
  \boxed{\frac{218}{3}-\frac{155}{2}
       =-\frac{29}{6}\approx-4.83\text{ years}.}
\]

\textbf{(b)} There are $80+20=100$ smokers and $40+60=100$
non-smokers among $200$ participants. Use common weights $1/2,1/2$:
\[
  \boxed{\text{Adjusted ATE estimate}
    =\frac12(68-70)+\frac12(82-80)=0\text{ years}.}
\]
Equivalently, the standardized means are
$\tfrac12(68)+\tfrac12(82)=75$ for coffee and
$\tfrac12(70)+\tfrac12(80)=75$ for no coffee.

\textbf{(c)} Smokers constitute $80/120=2/3$ of the coffee group but only
$20/80=1/4$ of the no-coffee group. Smokers have lower mean lifespan in both
groups, so this different composition pulls the raw contrast downward.
Using the same smoking mix removes that composition difference.
The stratum-specific estimates are $68-70=-2$ years for smokers and
$82-80=+2$ years for non-smokers. They cancel under the target's equal weights;
neither stratum estimate is zero. Causal interpretation depends on the stated
assumptions; these observational summaries alone do not establish them.
\end{answer}

\newpage
\subsection*{Problem 4: Rolling until the first six \hfill [8 minutes]}
Roll a fair six-sided die independently until the first $6$ appears.
Let $\rv{n}$ be the \emph{total} number of rolls, including the final $6$.
Let $A$ be the event that at least one $1$ appears before that first $6$.
\begin{enumerate}
  \item Find $\E[\rv{n}]$.
  \item Find $\E[\rv{n}\mid A]$, the expected total number of rolls given
    that at least one $1$ appears before stopping.
  \item Find $\E[\rv{n}\mid A^c]$, the expected total number of rolls given
    that no $1$ appears before stopping.
\end{enumerate}
\workingspace{8cm}

\begin{answer}
\textbf{(a)} Each roll is a $6$ with probability $1/6$, independently.
Thus $\rv{n}$ is geometric on $\{1,2,\ldots\}$ with success probability
$1/6$, and $\boxed{\E[\rv{n}]=6}$.

\paragraph{A useful intermediate stopping time.}
Let $\rv{w}$ be the number of rolls until the first $1$ or $6$, including
that roll. Its success probability is $2/6=1/3$, so $\E[\rv{w}]=3$.
The first such face is $1$ on $A$ and $6$ on $A^c$.
For $k=1,2,\ldots$,
\[
  \Pbb(\rv{w}=k,A)=\Pbb(\rv{w}=k,A^c)
  =\left(\frac23\right)^{k-1}\frac16
  =\frac12\Pbb(\rv{w}=k).
\]
Summing over $k$ gives $\Pbb(A)=\Pbb(A^c)=1/2$; the displayed
factorization also shows that $\rv{w}$ is independent of $A$. Hence
\[
  \E[\rv{w}\mid A]=\E[\rv{w}\mid A^c]=3.
\]

\textbf{(b)} On $A$, the first $1$ or $6$ is a $1$. After that roll,
the future rolls are independent of the past, so the additional waiting
time for a $6$ has mean $6$. The total includes both stages:
\[
  \boxed{\E[\rv{n}\mid A]=3+6=9.}
\]
The first stage includes the first $1$; the second stage includes the final $6$.

\textbf{(c)} On $A^c$, the first $1$ or $6$ is already a $6$, so
$\rv{n}=\rv{w}$ and
\[
  \boxed{\E[\rv{n}\mid A^c]=3.}
\]
As a check, iterated expectation gives
\[
  \E[\rv{n}]
  =\Pbb(A)\E[\rv{n}\mid A]+\Pbb(A^c)\E[\rv{n}\mid A^c]
  =\frac12(9)+\frac12(3)=6.
\]

\paragraph{Why the answer in (c) is not $5$.}
Conditioning on the entire stopped sequence containing no $1$ favors
shorter sequences. It does not make the remaining five faces equally likely
on each roll. In fact,
$\Pbb(\rv{n}=1\mid A^c)=(1/6)/(1/2)=1/3$, not $1/5$.
\end{answer}

\newpage
\subsection*{Problem 5 (backup): Splitting a Poisson total}
\noindent\textit{Optional if time allows. Suggested working time: 5--7 minutes.}
\medskip

Let $\rv{x}$ and $\rv{y}$ be independent
$\operatorname{Poisson}(\lambda)$ random variables, with $\lambda>0$,
and define $\rv{s}=\rv{x}+\rv{y}$.
\begin{enumerate}
  \item For any integer $n\geq0$, find
    $\mu_{\rv{x}\mid\rv{s}}(n)=\E[\rv{x}\mid\rv{s}=n]$.
    \textit{Hint:} use symmetry and $\rv{x}+\rv{y}=n$ under the condition.
  \item Write the random conditional mean
    $\mu_{\rv{x}\mid\rv{s}}(\rv{s})$ and use iterated expectation
    to recover $\E[\rv{x}]$.
\end{enumerate}
\workingspace{7cm}

\begin{answer}
The joint law of $\rv{x}$ and $\rv{y}$ is unchanged when they are exchanged,
and the event $\rv{x}+\rv{y}=n$ has the same symmetry. Thus
\[
  \E[\rv{x}\mid\rv{s}=n]=\E[\rv{y}\mid\rv{s}=n].
\]
Conditional linearity gives
\[
  \E[\rv{x}\mid\rv{s}=n]+\E[\rv{y}\mid\rv{s}=n]
  =\E[\rv{x}+\rv{y}\mid\rv{s}=n]=n.
\]
Each conditional mean is therefore $\boxed{n/2}$, including when $n=0$.
The corresponding random variable is
\[
  \boxed{\mu_{\rv{x}\mid\rv{s}}(\rv{s})=\frac{\rv{s}}2.}
\]
Since $\E[\rv{s}]=\E[\rv{x}]+\E[\rv{y}]=2\lambda$, the tower law gives
\[
  \E[\rv{x}]=\E\!\left[\frac{\rv{s}}2\right]=\lambda.
\]
As a distributional check,
$\rv{x}\mid\rv{s}=n\sim\operatorname{Binomial}(n,1/2)$,
whose mean is $n/2$.
\end{answer}

\vfill
\noindent\textit{Source:} adapted from Problem 3(b) of
\texttt{wyt/Recitation 4.pdf} in the Fall 2025 archive.

\ifsolutions
\newpage
\section*{Instructor pacing guide (75 minutes)}
\begin{center}
\begin{tabular}{@{}lll@{}}
\toprule
Time & Activity & Suggested use\\
\midrule
0--12 min & Conditional mean & Function vs. random variable; tower law\\
12--25 min & Average treatment effect & Consistency, randomization, adjustment\\
25--34 min & Problem 1: Exercise 7.11 & 5 min individual work, 4 min discussion\\
34--46 min & Problem 2 & 8 min individual work, 4 min discussion\\
46--62 min & Problem 3: Exercise 7.17 & 10 min individual work, 6 min discussion\\
62--74 min & Problem 4: Die rolls & 8 min individual work, 4 min discussion\\
74--75 min & Wrap-up & Check the tower law in Problem 4\\
\bottomrule
\end{tabular}
\end{center}

\paragraph{Review priorities.}
Start with the distinction between
$\mu_{\rv{y}\mid\rv{x}}(x)$ and
$\mu_{\rv{y}\mid\rv{x}}(\rv{x})$.
Explain iterated expectation as a weighted average of group means.
Then use the same idea to derive adjustment for confounding: first average
within each stratum under each treatment, then use common population weights.
Connect potential outcomes and consistency to Recitation 2.

\paragraph{Optional regression connection.}
If students ask why the conditional mean is useful for prediction, mention
that it minimizes mean squared error: when second moments are finite,
$h(x)=\mu_{\rv{y}\mid\rv{x}}(x)$ minimizes
$\E[(\rv{y}-h(\rv{x}))^2]$ over predictors based on $\rv{x}$.
The proof is outside today's core practice.

\paragraph{Common mistakes to catch.}
\begin{itemize}[leftmargin=1.5em,itemsep=0.45em]
  \item Using a joint distribution in a conditional-mean calculation without
    dividing by the conditioning probability.
  \item Averaging group means equally when group probabilities differ.
  \item Treating the observed coffee contrast as a causal effect without
    an independence assumption.
  \item Adjusting the coffee mean with coffee-group weights and the no-coffee
    mean with no-coffee-group weights. That reproduces the raw contrast.
  \item Interpreting zero overall adjusted effect as zero effect in every
    stratum. Here the estimates are $-2$ and $+2$ years.
  \item In Problem 4, treating ``no $1$ before the first $6$'' as independent
    rolls of a die with five equally likely faces.
\end{itemize}

\paragraph{Closing check.}
Ask: ``Which distribution supplies the weights in the tower law? Which
weights do we use to compare two treatments in the same target population?''
The added die problem makes the schedule tighter. If needed, provide the
conditional densities in Problem 2 so students can focus on its means.
Use Problem 5 only if the class is ahead. For Problem 4, a useful hint is to
stop first at either $1$ or $6$.

\paragraph{Sources and adaptations.}
The review follows Sections 7.8--7.9 of Carlos Fernandez-Granda's
\emph{Probability and Statistics for Data Science}.
Problems 1 and 3 rephrase Exercises 7.11 and 7.17, retaining their numerical
data; explanatory subparts have been added.
Problem 2 reuses the uniform--Bernoulli model from
\texttt{mine/r4/main.tex} in the Fall 2025 archive and adds conditional means
and iterated expectation. Problem 4 is the requested die-stopping exercise.
Problem 5 reuses the Poisson conditional-mean
question from \texttt{wyt/Recitation 4.pdf}, Problem 3(b), in that archive
(the original PDF is labeled Spring 2024--25).
\fi

\end{document}
